% ECONOMICS_PROBLEM: {"id": "D60-518", "title": "Welfare from inconsistent choices", "jel_code": "D60", "source_id": "OP-0273"}
% !TeX program = xelatex
\documentclass[11pt,a4paper]{article}
\usepackage[margin=23mm]{geometry}
\usepackage{fontspec}
\setmainfont{Latin Modern Roman}
\usepackage{amsmath,amssymb,mathtools}
\usepackage{xcolor,enumitem,booktabs,longtable,array,xurl,hyperref}
\definecolor{muted}{HTML}{53636C}
\hypersetup{colorlinks=true,urlcolor=blue,linkcolor=blue}
\setlength{\parindent}{0pt}
\setlength{\parskip}{5pt}
\newcommand{\R}{\mathbb R}
\newcommand{\E}{\mathbb E}
\newcommand{\N}{\mathbb N}
\newcommand{\ind}{\mathbf 1}
\newcommand{\Var}{\operatorname{Var}}
\newcommand{\Cov}{\operatorname{Cov}}
\newcommand{\supp}{\operatorname{supp}}
\DeclareMathOperator*{\argmin}{arg\,min}
\DeclareMathOperator*{\argmax}{arg\,max}
\newcommand{\entrymeta}[3]{{\sffamily\small\color{muted}\textbf{\texttt{#1}}\quad|\quad #2\par #3\par}}
\newcommand{\Problemref}[1]{\texttt{#1}}
\newcommand{\JELref}[2]{\texttt{#2} (JEL \texttt{#1})}
\begin{document}
\section*{Welfare from inconsistent choices}
\textbf{Problem ID:} \texttt{D60-518}\quad \textbf{JEL:} \texttt{D60}\par
% BEGIN ECONOMICS STATEMENT
\label{problem:OP-0273}
{\sffamily\small\color{muted}Original subject group: Economic theory, equilibrium and social choice\par}
\label{definitions:problem:30007554}
\textbf{Audit classification:} Published research agenda. Verified welfare discussion in §6.4 and enriched-data agenda in §7. The robust relation is editorial, not a unique welfare criterion endorsed by the survey.
\subsection*{Definitions and precise research target}
\textbf{Model and research target.} Let \(X\) be a finite set of economic alternatives, \(B\subseteq X\) a feasible menu, and \(f\in F\) an observed frame. For each menu and frame, observe choice probabilities \(P(x\mid B,f)\) and optional process measurements \(L\), such as response time or attention. Fix a publicly specified class \(\mathcal M\) of behavioral models. A model \(m\in\mathcal M\) consists of a latent welfare ordering \(\succeq_m\), a choice procedure, and a distribution of \(L\). State explicitly the normative assumption connecting that ordering to welfare; inconsistent choices alone cannot supply it. Let \(\mathcal M(P,L)\) contain all models fitting the population observables. Define robust welfare comparison by
\[
 x\succeq^R y\quad\Longleftrightarrow\quad x\succeq_m y\text{ for every }m\in\mathcal M(P,L).
\]
Determine the sharp identified relation \(\succeq^R\), which additional experimentally manipulated frames or process measurements enlarge it, and which comparisons remain intrinsically ambiguous within the maintained model class. Construct uncertainty regions accounting for sampling error and evaluate policy rankings by their robust welfare implications. A successful solution must report how conclusions change when competing attention or preference-construction models are admitted.

\emph{Definitions and maintained assumptions (editorial unless a source definition is named).} Menus are nonempty members of a declared \(\mathcal B\subseteq2^X\); \(P(x\mid B,f)=0\) for \(x\notin B\), and probabilities sum to one. A welfare ordering is a complete transitive binary relation on \(X\), linked to well-being by an explicitly defended normative assumption. Specify a measurable process space \(\mathcal L\), the population joint laws \(P^{B,f}\) of \((x,L)\), and a model-generated law \(P_m^{B,f}\). Define \(\mathcal M(P)=\{m:P_m^{B,f}=P^{B,f}\ \forall(B,f)\}\), and require it to be nonempty before taking its intersection of welfare relations. Otherwise universal quantification over the empty class would misleadingly rank every pair. A comparison is sharp when it holds in every fitting model and each excluded comparison has a fitting countermodel. With sampling, replace exact laws by a simultaneous \((1-\delta)\) confidence region for laws before forming the robust relation. Experiments shrink the fitting class only when their exclusion restrictions preserve the normative welfare ordering.
\subsection*{Variants and assumption changes}
\begin{enumerate}
\item \textbf{One procedure (editorial).} Specify one attention mechanism and its admissible parameters. Characterize the identified welfare relation and the smallest additional menu family needed to identify each remaining pair.
\item \textbf{Model union (editorial).} Allow attention, satisficing and context-dependent preference construction simultaneously. Compute comparisons surviving their union and quantify how results change under an explicit alternative normative interpretation of constructed preferences.
\end{enumerate}
\subsection*{References and source audit}
\begin{enumerate}
\item §6.4 pp.51–52; §7 pp.53–54 in February 2023 manuscript. Supports the stated research area; exact displayed specification is editorial. DOI publication is2024; full inspected manuscript is February2023. \url{https://geoffroydeclippel.net/Working%20Papers%20PDFs/JEL.pdf}
\end{enumerate}
\begin{enumerate}\raggedright
\item\hypertarget{definitions:source:30007554:1}{} Geoffroy de Clippel; Kareen Rozen. 2024. \emph{Bounded Rationality in Choice Theory: A Survey}. Published JEL 62(3),995–1039 (2024); inspected February2023 author manuscript §§6.4–7, printed pp.51–54.
\url{https://geoffroydeclippel.net/Working%20Papers%20PDFs/JEL.pdf}.
DOI: \url{https://doi.org/10.1257/jel.20231592}.
\end{enumerate}

\subsection*{Common conventions: Source mathematical conventions}

These conventions apply to each entry unless explicitly replaced.

\textbf{Models and identification.} A primitive model \(m\) specifies all probability laws, feasible choices, information, equilibrium and measurement rules used by its target. The admissible class \(\mathcal M\) is fixed independently of outcomes. If \(P_O(m)\) is its observable probability law and \(\tau(m)\) the target, its identified set at observable law \(P\) is
\[
 \mathcal I_\tau(P)=\{\tau(m):m\in\mathcal M,\ P_O(m)=P\}.
\]
Point identification means that this set is a singleton. An empty set signals incompatibility of data and assumptions. Coordinate bounds need not describe the joint set. Bounds are sharp only if every retained target is attainable or arbitrarily approachable by compatible models. Finite bounds require explicit restrictions on unobserved quantities; unrestricted confounding, hidden wealth, extrapolation or arbitrary equilibrium selection can yield uninformative sets.

\textbf{Causal interventions.} A potential outcome \(Y_i(p)\) is the outcome under a complete policy regime \(p\), including timing, financing, information and market spillovers. The target population and units are fixed before treatment. With interference, use the complete assignment vector or a justified exposure mapping; individual no-interference notation alone cannot identify an economy-wide intervention. A difference of mean potential outcomes does not require the cross-world joint distribution. Conditioning on post-treatment migration, survival, enrollment or employment changes the estimand and needs separate justification.

\textbf{Statistical statements.} An estimator, confidence set, forecast or test specifies a sampling law, model class and loss/coverage criterion. Uniform \((1-\alpha)\)-coverage means \(\inf_{m\in\mathcal M}\Pr_m\{\tau(m)\in C_n\}\geq1-\alpha\), or an explicitly asymptotic version. The whole parameter space trivially has coverage, so informativeness requires a stated width, risk or power objective. A forecasting improvement is relative to a named benchmark, information set and evaluation population.

\textbf{Equilibrium and optimization.} An equilibrium comprises feasible plans, optimality, beliefs where needed, budget constraints and all market/resource clearing. The primitives or a stated benchmark define each correspondence \(\mathcal E(p;m)\). Nonemptiness is assumed or proved, never inferred just from compact policy sets. A selected-equilibrium target uses a declared selection rule; a robust target quantifies over all permitted equilibria. Use suprema or infima unless attainment is established. An \(\varepsilon\)-optimal policy is feasible and within \(\varepsilon\) of the finite optimum value. Report an empty feasible set separately.

\textbf{Welfare and accounting.} Normative utility, interpersonal normalization, social weights, numeraire, discounting and the use of public receipts are primitives. Behavioral utility need not coincide with the welfare criterion. Household welfare includes ownership income and rents once; adding profits again to compensating variation generally double-counts them. A monetary equivalent is defined by an explicit transfer and price convention, with monotonicity and range conditions. Average effects, mechanism contributions, transfers and welfare losses are different targets. Infinite sums require convergence and dynamic feasible plans require terminal or no-Ponzi conditions.

\textbf{Source versions.} A working-paper date, revision date and journal publication year are distinguished. A DOI for a journal version is not automatically the DOI of an earlier manuscript. Locators refer to the stated version and printed pages where specified. A metadata-only check does not verify a claim about a full-text conclusion. Every entry's source subsection narrows the provenance claim accordingly.
% END ECONOMICS STATEMENT
\end{document}
